%%%%%
%
%% Dans ce fichier, les �ventuels accents sont cod�s
%% en ISO-8859-15 (aka latin9)
%% La ligne suivante donne trois exemples 
%% e accent aigu : �, e accent grave : �, e-dans-l'o : �
\section{Calibration par corr\'elation}

\itemstyle{Op\'erateur~:} \operfmt{Correl}
\index{Correl}

%%%%%OBJET%%%%
\itemstyle{Objet~:}
Calibration d'un vecteur1 avec le vecteur2 par r\'egression lin\'eaire selon
le principe des moindres carr\'es ordinaires appliqu\'e de proche en proche sur un intervalle
d\'efini par l'argument fenetre (en secondes).


%%%%%ARGUMENTS%%%%
\bigskip
%%%%%ENTREES%%%%
\begin{tabarg}
\multicolumn{3}{l}{\titretab {Entr\'ees}}\\
\hline
$Vecteur1$ & Vecteur, \texttt{float} & Vecteur \`a calibrer\\
$Vecteur2$ & Vecteur, \texttt{float} & Vecteur r\'ef\'erence pour la calibration\\
$fr$ & Coeff., \texttt{float} & Fr\'equence de l'\'echantillon\\
$fenetre$ & Coeff.,\texttt{float} & intervalle de calcul en secondes\\
 
\end{tabarg}

\medskip
%%%%%SORTIES%%%%
\begin{tabarg}
\multicolumn{3}{l}{\titretab {Sorties}}\\
\hline
$Vect$ & Vecteur, \texttt{float} & Vecteur r\'esultat\\
$Vect\_a$ & Vecteur, \texttt{float} & Vecteur coefficient a de la droite de r\'egression\\
$Vect\_b$ & Vecteur, \texttt{float} & Vecteur coefficient b de la droite de r\'egression\\
$Vect\_r2$ & Vecteur, \texttt{float} & Vecteur coefficient de corr\'elation lin\'eaire au carr\'e\\
\end{tabarg}


%%%%%FORMULE%%%%
\itemstyle{Formules}
%%Correl

Moyenne :
\[
\overline{X} = \frac{1}{n} \sum_{i = 1}^{i = n} X[i]
\]
Variance :
\[
Var(X) = \frac{1}{n}  \sum_{i = 1}^{i = n} {(X_i - \overline{X})}^2
\]
Covariance :
\[
Cov(X,Y) = \frac{1}{n}  \sum_{i = 1}^{i = n} (X_i - \overline{X})(Y_i - \overline{Y})
\]
Coefficients de la droite de r\'egression y = ax + b :
\[
a = \frac{Cov(X,Y)}{Var(X)}
\]
\[
b = \overline{Y} - a\overline{X}
\]
Coefficient de corr\'elation lin\'eaire :
\[
R = \frac{Cov(X,Y)}{\sqrt{Var(X)Var(Y)}}
\]


%%%%%COMMENTAIRE%%%%
\itemstyle{Commentaire~:}
Les 2 vecteurs doivent \^etre synchrones et � la m\^eme fr\'equence d'\'echantillonnage.

%%%%%AUTEUR%%%%
\itemstyle{Auteur~:}
CNRM/GMEI/TRAMM

%%%%%REFERENCE%%%%
